%------------------------------------------------------------------------------------
% ВВЕДЕНИЕ
%------------------------------------------------------------------------------------
\newpage
\begin{center}
  \textbf{\large ВВЕДЕНИЕ}
\end{center}
\addcontentsline{toc}{chapter}{ВВЕДЕНИЕ}

% TODO: Хочется написать про развитие телекоммуникаций
%В настоящий момент актуальной задачей развития телекоммуникаций является создание и развитие обьединенных сетей
%наземной сотовой связи (TN, Terrestrial Networks) и неназемной сотовой связи (NTN, Non Terrestrial Networks).
%Наземная сотовая связь состоит из вышек базовых станций, в то время как неназемная связь базируется в первую
%очередь на спутниках околоземной орибты (LEO), а также аэростатах, дронах и других стредствах (ТУТ НУЖНО УКАЗАТЬ ТЕРМИН ДЛЯ
%ТАКИХ СТРЕДСТВ из спеуификации 5G). Средства NTN позволяют разрешить многие задачи, с которыми столкнулись при разработке наземных систем связи,
%например, обеспечение полного покрытия сотовой (ПРОВЕРЬ ТЕРМИН СОТОВОЙ) связью всей поверхности Земли, обеспечение более
%качественной передачи данных на конкретных учатках, обслуживание потребностей интернета вещей, цифровых двойников и тд.

В настоящий момент актуальной задачей развития телекоммуникаций является создание и развитие объединенных сетей
наземной сотовой связи (TN, Terrestrial Networks) и неназемной сотовой связи (NTN, Non-Terrestrial Networks).
Наземная сотовая связь состоит из вышек базовых станций, в то время как неназемная связь базируется в первую очередь
на спутниках низкой околоземной орбиты (LEO, Low Earth Orbit), а также на высотных платформах (HAPS, High Altitude Platform Stations)
— аэростатах, дирижаблях и беспилотных летательных аппаратах (UAVs, Unmanned Aerial Vehicles) [\ref{Figaro2026}]. Средства NTN позволяют разрешить многие задачи,
с которыми столкнулись при разработке наземных систем связи. Например, обеспечение полного глобального покрытия 
связью всей поверхности Земли (включая Арктику, океаны и удаленные территории), обеспечение более качественной передачи
данных на конкретных участках (городская застройка, транспортные коридоры), обслуживание потребностей интернета вещей (IoT),
цифровых двойников и систем мониторинга [\ref{Rinaldi2020}].

% TODO: Далее мы пишем про выгоды спутниковых группировок
%NTN будет полезен в случае городской застройки, древенской, удаленной местности (Камчатка и тд)
%самолетах и кораблях.
NTN будет особенно полезен в условиях плотной городской застройки (для разгрузки наземных вышек),
в сельской и удаленной местности (например, Камчатка, Сибирь, Дальний Восток),
а также для обеспечения связью подвижных объектов: пассажирских самолетов, морских судов и поездов.
Внедрение спутникового сегмента в архитектуру сетей 5G позволяет решить проблему цифрового неравенства,
обеспечивая широкополосным доступом территории с низкой плотностью населения,
где развертывание наземной инфраструктуры экономически нецелесообразно. 
Кроме того, космическая связь выступает в качестве физически распределенного резервного канала, 
устойчивого к наземным техногенным авариям и стихийным бедствиям, что критически повышает отказоустойчивость 
национальной телекоммуникационной инфраструктуры, особенно в условиях чрезвычайных ситуаций.

% TODO: Далее можно будет вставить небольшой обзор того, что сделано в 5G NR NTN
% TODO: И вообще описать, что тема новая, работы ведутся во всю
%Современные системы спутниковой мобильной связи базируются на техноголии стандарта 3gpp 5G NR (New Radio, МОЖНО ВАЛЬНУТЬ ССЫЛКУ),
%на основании этого стандарта был разработан стандарт спутниковой связи 5G NR NTN (ССЫЛКА НА НЕГО).

%Собщество 3gpp вело разработки и на данный момент (Rel 19) было определено это и это, такие то пути пройдены,
%такие выводы сделаны. (ССЫЛКА НА КАКОЙ НИБУДЬ ОБЗОР).

%Современные системы спутниковой мобильной связи базируются на технологии стандарта 3GPP 5G NR (New Radio).
%На основании этого стандарта был разработан специализированный стандарт спутниковой связи 5G NR NTN.
%Сообщество 3GPP вело активные разработки в этом направлении, и на данный момент (вплоть до Release 18)
%определены ключевые архитектурные решения. В частности, в Release 17 была заложена основа для NTN 
%с «прозрачной» (ретрансляционной) полезной нагрузкой спутников и поддержка сценариев IoT-NTN.
%В Release 18 (5G-Advanced) были внесены улучшения для эфемерид, компенсации доплеровского сдвига на абонентском оборудовании (UE),
%поддержка работы в Ka-диапазоне и повышение мобильности/непрерывности обслуживания.

Современные системы спутниковой мобильной связи базируются на технологии стандарта 3GPP 5G NR (New Radio).
На основании этого стандарта был разработан специализированный стандарт спутниковой связи 5G NR NTN[\ref{TR38811}].
Сообщество 3GPP вело активные разработки в этом направлении, и на данный момент (вплоть до Release 20)
определены ключевые архитектурные решения. В Release 17 была заложена основа для NTN 
с «прозрачной» (ретрансляционной) полезной нагрузкой спутников и поддержка сценариев IoT-NTN.
В Release 18 (5G-Advanced) были внесены улучшения для эфемерид, компенсации доплеровского сдвига на абонентском оборудовании (UE),
поддержка работы в Ka-диапазоне и повышение мобильности/непрерывности обслуживания.
Release 19 вводит регенеративную полезную нагрузки (базовая станция на борту), а
Release 20  закладывает основы для 6G-NTN.

% TODO: Поэтому было бы неплохо сделать такую систему и у нас
%В данный момент актуальной темой является создание собственной сети спутниковой
%связи, которая поможет развитию и тд и тп. Данными задачами занимаются многие компании 
%(Бюро 1440, АО Решентнев и тд.). Поэтому важно вести разработки и исследования в этой области.

В России созданием национальных спутниковых сетей занимаются ГК «Роскосмос»,
АО «Решетнёв» и «Бюро 1440». АО «Решетнёв» разрабатывает систему «Марафон IoT»
(132 спутника для передачи данных с датчиков).
«Бюро 1440» реализует проект низкоорбитальной группировки «Рассвет» с поддержкой 5G NTN.
ГК «Роскосмос» координирует создание среднеорбитальной системы «Скиф» для широкополосного доступа.
Развитие этих проектов необходимо для технологического суверенитета РФ в области спутниковой связи.

% TODO: Подходим к теме исследования
%Современные технологии космической связи требуют высокой надежности и
%эффективности передачи данных, что делает актуальным исследование
%различных факторов, влияющих на качество сигналов. Одним из таких факторов
%является эффект Доплера, который возникает в результате относительного
%движения источника и приемника сигналов. В современных системах передачи
%данных [\ref{steputin20216g_v1}],
%использующих ортогональное частотное деление (OFDM,
%Orthogonal Frequency Division Multiplexing), влияние эффекта Доплера может
%существенно сказаться на характеристиках принимаемого сигнала, таких как
%уровень искажений, скорость передачи данных и устойчивость к помехам.

Современные технологии космической связи требуют высокой надежности и эффективности передачи данных,
что делает актуальным исследование различных факторов, влияющих на качество сигналов.
Одним из таких факторов является эффект Доплера, который возникает в результате относительного
движения источника и приемника сигналов. В системах связи с ортогональным частотным разделением каналов [\ref{steputin20216g_v1}]
(OFDM, Orthogonal Frequency Division Multiplexing), влияние эффекта Доплера может существенно сказаться
на характеристиках принимаемого сигнала, таких как уровень искажений, скорость передачи данных и устойчивость к помехам.

%Модуляция OFDM является одной из наиболее перспективных технологий для
%реализации высокоскоростной передачи информации в условиях ограниченной
%полосы пропускания и каналов связи c замираниями. Благодаря своей способности
%эффективно противостоять межсимвольной интерференции и многолучевому
%распространению, модуляция OFDM широко применяется в системах связи.

%Однако при
%наличии значительных скоростей движения приемников или передатчиков,
%таких как спутники или космические аппараты, эффект Доплера приводит к
%смещению частот поднесущих,
% что в свою очередь вызывает проблемы с
%синхронизацией и декодированием информации.

%Наряду со смещением частот, для борьбы с которым есть множество методов,
%возникает эффект масштабирования спектра, который приводит к возникновению
%интерференции между поднесущими (ICI, Inter Carrier Interference).
%Так как основной принцип OFDM модуляции заключается в ортогональности
%поднесущих, данный метод особенно чувствителен к влиянию ICI.

Модуляция OFDM является одной из наиболее перспективных технологий для реализации высокоскоростной передачи
информации в условиях ограниченной полосы пропускания и каналов связи с замираниями.
Благодаря своей способности эффективно противостоять межсимвольной интерференции (ISI)
и многолучевому распространению, OFDM широко применяется в системах связи, включая стандарт 5G NR.
Однако при наличии значительных скоростей движения приемников или передатчиков,
характерных для спутников на низкой околоземной орбите,
эффект Доплера приводит не только к смещению несущей частоты,
но и к масштабированию спектра. Это явление вызывает потерю ортогональности между поднесущими,
что в свою очередь порождает межканальную интерференцию (ICI, Inter-Carrier Interference).

Наряду со смещением частот, для борьбы с которым существует множество классических методов
(пилот-сигналы, автоподстройка частоты), эффект масштабирования спектра является более сложной проблемой.
Так как основной принцип OFDM заключается в строгой ортогональности поднесущих, данный метод модуляции 
особенно чувствителен к ICI, которая приводит резкому снижению помехоустойчивости при росте отношения сигнал/шум.

% TODO: Тут можно вставить работы, исследующие методы борьбы с Допплером
% TODO: Мини обзор того, что было сделано
%Многие исследования были направлены на изучение влияния доплеровского сдвига и
%мастабирования спектра на производительность системы связи и на методы компенсации
%данного эффекта. (ТУТ НУЖНО ДОБАВИТЬ ССЫЛКИ НА ИССЛЕДОВАНИЯ).

Многие исследования были направлены на изучение влияния доплеровского сдвига и масштабирования спектра
на производительность систем связи, а также на разработку методов компенсации данных эффектов [\ref{Figaro2023}, \ref{Asciolla2025}, \ref{Gannon2022}].
В научной литературе предлагаются различные подходы: от временного окна (windowing)
и частотной коррекции до использования частичного БПФ (Partial FFT) и адаптивной фильтрации.
Однако большинство работ либо сфокусированы на наземных сценариях высокоскоростного транспорта (поезда до 500 км/ч),
либо рассматривают спутниковые каналы в общем виде, без привязки к жестким спецификациям
физического уровня 5G NR NTN. Поэтому важно провести исследование влияния эффекта доплеровского
масштабирования спектра на реальную систему связи, на примере стандарта 5G NR-NTN.

% TODO: Оглашение самой гипотезы
%Предполагается, что есть такие условия (парметры системы, например numerologies), не описанные в спецификации,
%которые приводят к существенному снижению эффективности системы (падение BER ~ -6 дБ от TN).
%Но есть способы, которые позволяют этот эффект невелированть (до -(1-2) дБ).

\textbf{Гипотеза исследования.} Предполагается, что существующие в спецификации 5G NR NTN наборы параметров,
а именно расстояние между поднесущими, по-разному чувствительны к эффекту масштабирования.
Существуют такие условия (скорости спутника, рабочая частота), которые приводят к существенному снижению эффективности системы
(падение BER до уровня, сравнимого с ухудшением на 6 дБ относительно наземного сценария).
%При этом предполагается, что применение адаптивных методов обработки сигнала 
%(например, использование менее чувствительных numerology или специальных схем коррекции на уровне приемника)
%позволяет нивелировать этот эффект до приемлемого уровня (ухудшение не более 1-2 дБ).


%------------------------------------------------------------------------------------
% ПОСТАНОВКА ЗАДАЧИ
%------------------------------------------------------------------------------------
\newpage
\begin{center}
  \textbf{\large ПОСТАНОВКА ЗАДАЧИ}
\end{center}
\addcontentsline{toc}{chapter}{ПОСТАНОВКА ЗАДАЧИ}

% TODO: Здесь нужно как можно более точно поставить свою цель
\textbf{Целью} работы является исследование влияния эффекта масштабирования спектра на 
произодительность системы спустниковой связи на примере системы стандарта 5G NR NTN 
для различных спутниковых орбит.

% TODO: Некоторая методологическая сноска
%В данной работе будут рассмотрены орбиты типа LEO и MEO.
%Стандарт 5G-NR NTN поддерживает LEO, MEO, GEO, а также HEO орбиты

%В данной работе будут рассмотрены орбиты типов LEO (низкая опорная орбита, высота ~500-2000 км)
%и MEO (средняя орбита, ~8000-20000 км). Стандарт 5G NR NTN также поддерживает GEO (геостационарную)
%и HEO (высокоэллиптическую) орбиты, однако в силу специфики доплеровских искажений 
%(максимальная скорость изменения задержки и частоты характерна именно для LEO) 
%основное внимание в исследовании уделяется низкоорбитальным и среднеорбитальным сегментам.

% TODO: Здесь нужно уточнить задачи
\textbf{Задачи:}
\begin{enumerate}
    \item Провести обзор литературы;
    \item получить теоретическию оценку ICI;
    \item провести численное моделирование системы;
    \item оценить влияние эффекта масштабирования спектра на производительность системы
\end{enumerate}

